\documentclass[aspectratio=169,11pt]{beamer}
\usetheme{Madrid}
\usecolortheme{default}
\usepackage[T1]{fontenc}
\setbeamertemplate{navigation symbols}{}
\setbeamertemplate{footline}[frame number]
\usepackage{enumitem}
\usepackage{amsmath,amssymb}
\usepackage{booktabs}
\usepackage{array}

\title{Databricks ML Associate}
\subtitle{High-Yield Top 100 (Most Frequently Tested Topics)}
\author{Practice Exam Deck}
\date{}

\begin{document}
	
	\begin{frame}
		\titlepage
	\end{frame}
	
	\begin{frame}{How to Use This Deck}
		\begin{itemize}
			\item 100 high-yield questions (Q1--Q100) on the topics that appear most often on the Databricks Machine Learning Associate exam.
			\item Each question slide has four options (A--D); the next slide gives the correct answer and why the others are wrong.
			\item Topic order: Databricks ML platform and AutoML, MLflow, Unity Catalog model registry, Feature Engineering, data processing, tuning, evaluation, deployment.
		\end{itemize}
	\end{frame}
	
	\begin{frame}{Exam Domains Covered}
		\begin{itemize}
			\item \textbf{Databricks Machine Learning:} runtime, AutoML, MLflow, Unity Catalog registry, Feature Engineering client.
			\item \textbf{Data Processing:} EDA, imputation, encoding, Spark ML transformers, pandas API on Spark and UDFs.
			\item \textbf{Model Development:} Hyperopt, CrossValidator, metrics, bias/variance, leakage.
			\item \textbf{Model Deployment:} batch, streaming, real-time serving, monitoring.
		\end{itemize}
	\end{frame}
	
	% ---------------------------------------------------------------
	\begin{frame}{Question 1}
		Which cluster runtime comes preinstalled with MLflow, scikit-learn, XGBoost, and Hyperopt?
		\begin{enumerate}[label=\Alph*.]
			\item Photon only
			\item Standard Databricks Runtime
			\item Databricks Runtime for Machine Learning
			\item A SQL warehouse
		\end{enumerate}
	\end{frame}
	\begin{frame}{Answer 1}
		\textbf{Correct: C} --- ML Runtime bundles the common ML libraries and tools.\\
		\vspace{0.5em}
		\textbf{Wrong:}
		\begin{itemize}
			\item A: Photon is a faster query engine, not an ML library bundle.
			\item B: does not bundle the ML library set.
			\item D: runs SQL, not notebook ML libraries.
		\end{itemize}
	\end{frame}
	
	% ---------------------------------------------------------------
	\begin{frame}{Question 2}
		A data scientist trains a scikit-learn model on a 2 GB dataset. Which cluster setup is most appropriate?
		\begin{enumerate}[label=\Alph*.]
			\item A 16-worker multi-node cluster
			\item A SQL warehouse
			\item An autoscaling cluster with many workers
			\item A Single Node cluster
		\end{enumerate}
	\end{frame}
	\begin{frame}{Answer 2}
		\textbf{Correct: D} --- scikit-learn runs on one machine, so extra workers would sit idle.\\
		\vspace{0.5em}
		\textbf{Wrong:}
		\begin{itemize}
			\item A: scikit-learn does not distribute across workers.
			\item B: not meant for Python ML training.
			\item C: workers would still go unused by scikit-learn.
		\end{itemize}
	\end{frame}
	
	% ---------------------------------------------------------------
	\begin{frame}{Question 3}
		Which problem types does Databricks AutoML support?
		\begin{enumerate}[label=\Alph*.]
			\item Classification, regression, and clustering
			\item Classification, regression, and forecasting
			\item Image generation and reinforcement learning
			\item Clustering, anomaly detection, and recommendation
		\end{enumerate}
	\end{frame}
	\begin{frame}{Answer 3}
		\textbf{Correct: B} --- These are the three supported AutoML problem types.\\
		\vspace{0.5em}
		\textbf{Wrong:}
		\begin{itemize}
			\item A: clustering is not supported.
			\item C: not AutoML problem types.
			\item D: none of these are supported.
		\end{itemize}
	\end{frame}
	
	% ---------------------------------------------------------------
	\begin{frame}{Question 4}
		Why is Databricks AutoML called a \texttt{}glass box'' solution?
		\begin{enumerate}[label=\Alph*.]
			\item It only works through a no-code UI
			\item It automatically deploys the best model to production
			\item It generates editable notebooks with the source code for each trial
			\item It hides all trial code and returns only the best model
		\end{enumerate}
	\end{frame}
	\begin{frame}{Answer 4}
		\textbf{Correct: C} --- you can inspect, edit, and reuse the generated code.\\
		\vspace{0.5em}
		\textbf{Wrong:}
		\begin{itemize}
			\item A: it also has a Python API.
			\item B: AutoML does not auto-deploy.
			\item D: that is a black box.
		\end{itemize}
	\end{frame}
	
	% ---------------------------------------------------------------
	\begin{frame}{Question 5}
		What does AutoML produce that helps you understand data quality before modeling?
		\begin{enumerate}[label=\Alph*.]
			\item A serving endpoint
			\item A data exploration notebook
			\item A Feature Store table
			\item A Hyperopt Trials object
		\end{enumerate}
	\end{frame}
	\begin{frame}{Answer 5}
		\textbf{Correct: B} --- it profiles the dataset and flags issues such as nulls and skew.\\
		\vspace{0.5em}
		\textbf{Wrong:}
		\begin{itemize}
			\item A: deployment is a separate step.
			\item C: AutoML does not create one for you.
			\item D: this is not an AutoML output.
		\end{itemize}
	\end{frame}
	
	% ---------------------------------------------------------------
	\begin{frame}{Question 6}
		Which call starts an AutoML experiment to predict a numeric target such as price?
		\begin{enumerate}[label=\Alph*.]
			\item \texttt{databricks.automl.cluster(df)}
			\item \texttt{mlflow.automl.run(df)}
			\item \texttt{databricks.automl.regress(df, target\_col="price")}
			\item \texttt{databricks.automl.classify(df, target\_col="price")}
		\end{enumerate}
	\end{frame}
	\begin{frame}{Answer 6}
		\textbf{Correct: C} --- regress is for continuous numeric targets.\\
		\vspace{0.5em}
		\textbf{Wrong:}
		\begin{itemize}
			\item A: there is no cluster API.
			\item B: not a real API.
			\item D: classify is for categorical targets.
		\end{itemize}
	\end{frame}
	
	% ---------------------------------------------------------------
	\begin{frame}{Question 7}
		After \texttt{summary = databricks.automl.classify(...)}, how do you access the best model's details?
		\begin{enumerate}[label=\Alph*.]
			\item \texttt{summary.champion}
			\item \texttt{summary.deploy()}
			\item \texttt{summary.models[0]}
			\item \texttt{summary.best\_trial}
		\end{enumerate}
	\end{frame}
	\begin{frame}{Answer 7}
		\textbf{Correct: D} --- the summary object exposes best\_trial with its model path and notebook.\\
		\vspace{0.5em}
		\textbf{Wrong:}
		\begin{itemize}
			\item A: not an attribute.
			\item B: AutoML summaries do not deploy.
			\item C: not how the summary is organized.
		\end{itemize}
	\end{frame}
	
	% ---------------------------------------------------------------
	\begin{frame}{Question 8}
		Which Databricks feature syncs notebooks with a Git provider such as GitHub?
		\begin{enumerate}[label=\Alph*.]
			\item Workspace trash
			\item Git folders (Repos)
			\item Unity Catalog volumes
			\item DBFS mounts
		\end{enumerate}
	\end{frame}
	\begin{frame}{Answer 8}
		\textbf{Correct: B} --- Git folders provide version control inside the workspace.\\
		\vspace{0.5em}
		\textbf{Wrong:}
		\begin{itemize}
			\item A: only holds deleted objects.
			\item C: file storage, not Git.
			\item D: storage access, not version control.
		\end{itemize}
	\end{frame}
	
	% ---------------------------------------------------------------
	\begin{frame}{Question 9}
		What is the correct Unity Catalog naming pattern for a registered model?
		\begin{enumerate}[label=\Alph*.]
			\item \texttt{workspace.model\_name}
			\item \texttt{schema.model\_name}
			\item \texttt{catalog.schema.model\_name}
			\item \texttt{cluster.catalog.model\_name}
		\end{enumerate}
	\end{frame}
	\begin{frame}{Answer 9}
		\textbf{Correct: C} --- UC uses a three-level namespace.\\
		\vspace{0.5em}
		\textbf{Wrong:}
		\begin{itemize}
			\item A: no such pattern.
			\item B: missing the catalog level.
			\item D: clusters are not part of the name.
		\end{itemize}
	\end{frame}
	
	% ---------------------------------------------------------------
	\begin{frame}{Question 10}
		A notebook needs a library that is not on the cluster. Which command installs it only for the current notebook session?
		\begin{enumerate}[label=\Alph*.]
			\item Editing the cluster policy
			\item \texttt{spark.conf.set()}
			\item \texttt{dbutils.fs.cp()}
			\item \texttt{\%pip install <package>}
		\end{enumerate}
	\end{frame}
	\begin{frame}{Answer 10}
		\textbf{Correct: D} --- notebook-scoped installs affect only that notebook's environment.\\
		\vspace{0.5em}
		\textbf{Wrong:}
		\begin{itemize}
			\item A: this does not install packages.
			\item B: sets Spark config, not libraries.
			\item C: copies files.
		\end{itemize}
	\end{frame}
	
	% ---------------------------------------------------------------
	\begin{frame}{Question 11}
		How should an API token be used in a notebook without exposing it in code?
		\begin{enumerate}[label=\Alph*.]
			\item Store it in a markdown cell
			\item Store it in a secret scope and read it with \texttt{dbutils.secrets.get()}
			\item Hard-code it in a variable
			\item Put it in a notebook widget default
		\end{enumerate}
	\end{frame}
	\begin{frame}{Answer 11}
		\textbf{Correct: B} --- secrets are redacted and access-controlled.\\
		\vspace{0.5em}
		\textbf{Wrong:}
		\begin{itemize}
			\item A: plain text and visible.
			\item C: exposes the credential.
			\item D: widgets are visible.
		\end{itemize}
	\end{frame}
	
	% ---------------------------------------------------------------
	\begin{frame}{Question 12}
		Where should you store unstructured files (such as images) under Unity Catalog governance?
		\begin{enumerate}[label=\Alph*.]
			\item A feature table
			\item A Delta table column of file names only
			\item The cluster's local disk
			\item A Unity Catalog volume
		\end{enumerate}
	\end{frame}
	\begin{frame}{Answer 12}
		\textbf{Correct: D} --- volumes govern non-tabular files.\\
		\vspace{0.5em}
		\textbf{Wrong:}
		\begin{itemize}
			\item A: for features, not files.
			\item B: does not store or govern the files.
			\item C: ephemeral and ungoverned.
		\end{itemize}
	\end{frame}
	
	% ---------------------------------------------------------------
	\begin{frame}{Question 13}
		A team wants a notebook to retrain a model every night at lower cost. What should they use?
		\begin{enumerate}[label=\Alph*.]
			\item A Git folder sync
			\item A SQL alert
			\item A Databricks Workflows job with a schedule on a job cluster
			\item Run manually on an all-purpose cluster
		\end{enumerate}
	\end{frame}
	\begin{frame}{Answer 13}
		\textbf{Correct: C} --- job clusters start for the run and terminate afterward.\\
		\vspace{0.5em}
		\textbf{Wrong:}
		\begin{itemize}
			\item A: does not run notebooks.
			\item B: alerts do not retrain models.
			\item D: not automated and costs more.
		\end{itemize}
	\end{frame}
	
	% ---------------------------------------------------------------
	\begin{frame}{Question 14}
		How can a team reproduce the exact training data used for an earlier model version in Delta Lake?
		\begin{enumerate}[label=\Alph*.]
			\item Run \texttt{OPTIMIZE}
			\item Re-run \texttt{df.cache()}
			\item Use Delta time travel with \texttt{versionAsOf} or \texttt{timestampAsOf}
			\item Run \texttt{VACUUM} with 0 hours
		\end{enumerate}
	\end{frame}
	\begin{frame}{Answer 14}
		\textbf{Correct: C} --- time travel reads earlier table versions.\\
		\vspace{0.5em}
		\textbf{Wrong:}
		\begin{itemize}
			\item A: compacts files.
			\item B: caching is not versioning.
			\item D: deletes old files and breaks time travel.
		\end{itemize}
	\end{frame}
	
	% ---------------------------------------------------------------
	\begin{frame}{Question 15}
		What is an MLflow experiment?
		\begin{enumerate}[label=\Alph*.]
			\item A single training execution
			\item A Delta table of features
			\item A deployed REST endpoint
			\item A named collection of runs for a particular task
		\end{enumerate}
	\end{frame}
	\begin{frame}{Answer 15}
		\textbf{Correct: D} --- runs are grouped under an experiment.\\
		\vspace{0.5em}
		\textbf{Wrong:}
		\begin{itemize}
			\item A: that is a run.
			\item B: that is a feature table.
			\item C: that is Model Serving.
		\end{itemize}
	\end{frame}
	
	% ---------------------------------------------------------------
	\begin{frame}{Question 16}
		Which code correctly creates a tracked run in a chosen experiment?
		\begin{enumerate}[label=\Alph*.]
			\item \texttt{mlflow.set\_experiment(path)} then \texttt{with mlflow.start\_run():}
			\item \texttt{mlflow.new\_run(path)}
			\item \texttt{mlflow.log\_run()} alone
			\item \texttt{mlflow.create\_experiment\_run()}
		\end{enumerate}
	\end{frame}
	\begin{frame}{Answer 16}
		\textbf{Correct: A} --- set the experiment, then open a run context.\\
		\vspace{0.5em}
		\textbf{Wrong:}
		\begin{itemize}
			\item B: not an API.
			\item C: does not start a run.
			\item D: not an API.
		\end{itemize}
	\end{frame}
	
	% ---------------------------------------------------------------
	\begin{frame}{Question 17}
		Which call logs a value that can be recorded repeatedly with a step, such as training loss per epoch?
		\begin{enumerate}[label=\Alph*.]
			\item \texttt{mlflow.log\_metric("loss", value, step=epoch)}
			\item \texttt{mlflow.set\_tag("loss", value)}
			\item \texttt{mlflow.log\_param("loss", value)}
			\item \texttt{mlflow.log\_artifact("loss")}
		\end{enumerate}
	\end{frame}
	\begin{frame}{Answer 17}
		\textbf{Correct: A} --- metrics are numeric and support steps.\\
		\vspace{0.5em}
		\textbf{Wrong:}
		\begin{itemize}
			\item B: tags are string metadata.
			\item C: params are fixed once set.
			\item D: logs files.
		\end{itemize}
	\end{frame}
	
	% ---------------------------------------------------------------
	\begin{frame}{Question 18}
		Which call saves a plot image file to a run?
		\begin{enumerate}[label=\Alph*.]
			\item \texttt{mlflow.set\_tag("plot.png")}
			\item \texttt{mlflow.log\_artifact("plot.png")}
			\item \texttt{mlflow.log\_metric("plot.png")}
			\item \texttt{mlflow.log\_param("plot.png")}
		\end{enumerate}
	\end{frame}
	\begin{frame}{Answer 18}
		\textbf{Correct: B} --- artifacts are files and directories attached to a run.\\
		\vspace{0.5em}
		\textbf{Wrong:}
		\begin{itemize}
			\item A: tags are text labels.
			\item C: metrics are numbers.
			\item D: params are small key-value pairs.
		\end{itemize}
	\end{frame}
	
	% ---------------------------------------------------------------
	\begin{frame}{Question 19}
		What is true about Databricks Autologging on ML Runtime clusters?
		\begin{enumerate}[label=\Alph*.]
			\item It works only for Spark ML
			\item It is enabled by default for supported libraries
			\item It only works after \texttt{mlflow.register\_model()}
			\item It must be set up with an init script
		\end{enumerate}
	\end{frame}
	\begin{frame}{Answer 19}
		\textbf{Correct: B} --- supported flavors log params, metrics, and models automatically.\\
		\vspace{0.5em}
		\textbf{Wrong:}
		\begin{itemize}
			\item A: many libraries are supported.
			\item C: unrelated.
			\item D: no init script is needed.
		\end{itemize}
	\end{frame}
	
	% ---------------------------------------------------------------
	\begin{frame}{Question 20}
		What does \texttt{mlflow.search\_runs()} return by default?
		\begin{enumerate}[label=\Alph*.]
			\item A registered model version
			\item A single best model
			\item A Spark Streaming query
			\item A pandas DataFrame of runs
		\end{enumerate}
	\end{frame}
	\begin{frame}{Answer 20}
		\textbf{Correct: D} --- it returns a DataFrame you can sort and filter.\\
		\vspace{0.5em}
		\textbf{Wrong:}
		\begin{itemize}
			\item A: it does not touch the registry.
			\item B: it returns run records.
			\item C: not what it returns.
		\end{itemize}
	\end{frame}
	
	% ---------------------------------------------------------------
	\begin{frame}{Question 21}
		Which call returns the run with the lowest RMSE in an experiment?
		\begin{enumerate}[label=\Alph*.]
			\item \texttt{mlflow.search\_runs(experiment\_ids=[id], order\_by=["metrics.rmse ASC"], max\_results=1)}
			\item \texttt{mlflow.get\_min\_run()}
			\item \texttt{mlflow.best\_run("rmse")}
			\item \texttt{mlflow.search\_runs(order\_by=["metrics.rmse DESC"], max\_results=1)}
		\end{enumerate}
	\end{frame}
	\begin{frame}{Answer 21}
		\textbf{Correct: A} --- sort by the metric ascending and keep one row.\\
		\vspace{0.5em}
		\textbf{Wrong:}
		\begin{itemize}
			\item B: not an API.
			\item C: not an API.
			\item D: this returns the highest RMSE.
		\end{itemize}
	\end{frame}
	
	% ---------------------------------------------------------------
	\begin{frame}{Question 22}
		What does an MLflow model signature describe?
		\begin{enumerate}[label=\Alph*.]
			\item The cluster used for training
			\item The input and output schema of the model
			\item The model's license
			\item The hyperparameter search space
		\end{enumerate}
	\end{frame}
	\begin{frame}{Answer 22}
		\textbf{Correct: B} --- it records expected column names and types.\\
		\vspace{0.5em}
		\textbf{Wrong:}
		\begin{itemize}
			\item A: not part of the signature.
			\item C: not part of the signature.
			\item D: not part of the signature.
		\end{itemize}
	\end{frame}
	
	% ---------------------------------------------------------------
	\begin{frame}{Question 23}
		Which function infers a model signature from data?
		\begin{enumerate}[label=\Alph*.]
			\item \texttt{mlflow.models.deploy(X)}
			\item \texttt{mlflow.sign(X)}
			\item \texttt{mlflow.schema.get(X)}
			\item \texttt{mlflow.models.infer\_signature(X, predictions)}
		\end{enumerate}
	\end{frame}
	\begin{frame}{Answer 23}
		\textbf{Correct: D} --- it derives schema from sample inputs and outputs.\\
		\vspace{0.5em}
		\textbf{Wrong:}
		\begin{itemize}
			\item A: deployment is unrelated.
			\item B: not an API.
			\item C: not an API.
		\end{itemize}
	\end{frame}
	
	% ---------------------------------------------------------------
	\begin{frame}{Question 24}
		Which MLflow flavor lets you load any logged model and call a generic \texttt{predict()}?
		\begin{enumerate}[label=\Alph*.]
			\item \texttt{mlflow.tracking}
			\item \texttt{python\_function} (pyfunc)
			\item \texttt{mlflow.projects}
			\item \texttt{mlflow.spark} only
		\end{enumerate}
	\end{frame}
	\begin{frame}{Answer 24}
		\textbf{Correct: B} --- pyfunc is the generic, framework-agnostic interface.\\
		\vspace{0.5em}
		\textbf{Wrong:}
		\begin{itemize}
			\item A: is the tracking API.
			\item C: packages code, not models.
			\item D: limited to Spark ML models.
		\end{itemize}
	\end{frame}
	
	% ---------------------------------------------------------------
	\begin{frame}{Question 25}
		How do you apply a logged MLflow model to a large Spark DataFrame?
		\begin{enumerate}[label=\Alph*.]
			\item \texttt{df.toPandas()} then predict on the driver
			\item \texttt{mlflow.pyfunc.spark\_udf(spark, model\_uri)}
			\item \texttt{model.collect()}
			\item \texttt{mlflow.score(df)}
		\end{enumerate}
	\end{frame}
	\begin{frame}{Answer 25}
		\textbf{Correct: B} --- it wraps the model as a Spark UDF for distributed scoring.\\
		\vspace{0.5em}
		\textbf{Wrong:}
		\begin{itemize}
			\item A: does not scale.
			\item C: models have no such method.
			\item D: not an API.
		\end{itemize}
	\end{frame}
	
	% ---------------------------------------------------------------
	\begin{frame}{Question 26}
		A team tunes hyperparameters and wants one parent run with a child run per trial. What should they use?
		\begin{enumerate}[label=\Alph*.]
			\item \texttt{mlflow.end\_run()} repeatedly
			\item Separate experiments per trial
			\item \texttt{mlflow.log\_param} for all trials in one run
			\item \texttt{mlflow.start\_run(nested=True)} inside a parent run
		\end{enumerate}
	\end{frame}
	\begin{frame}{Answer 26}
		\textbf{Correct: D} --- nested runs create child runs under a parent.\\
		\vspace{0.5em}
		\textbf{Wrong:}
		\begin{itemize}
			\item A: ends runs, not nesting them.
			\item B: loses the grouping.
			\item C: does not create child runs.
		\end{itemize}
	\end{frame}
	
	% ---------------------------------------------------------------
	\begin{frame}{Question 27}
		How can you visually compare many runs in the MLflow UI?
		\begin{enumerate}[label=\Alph*.]
			\item Select runs and click Compare to see parallel coordinates and scatter plots
			\item Open each run's artifact folder
			\item Check the cluster event log
			\item Use the Spark UI
		\end{enumerate}
	\end{frame}
	\begin{frame}{Answer 27}
		\textbf{Correct: A} --- the compare view helps relate parameters to metrics.\\
		\vspace{0.5em}
		\textbf{Wrong:}
		\begin{itemize}
			\item B: slow and not comparative.
			\item C: unrelated to run comparison.
			\item D: shows jobs, not runs.
		\end{itemize}
	\end{frame}
	
	% ---------------------------------------------------------------
	\begin{frame}{Question 28}
		How do you register a model at the time you log it?
		\begin{enumerate}[label=\Alph*.]
			\item Pass \texttt{registered\_model\_name=} to \texttt{mlflow.<flavor>.log\_model()}
			\item Pass \texttt{register=True} to \texttt{mlflow.start\_run()}
			\item Call \texttt{mlflow.sklearn.save\_model()}
			\item Use \texttt{mlflow.log\_param("register", True)}
		\end{enumerate}
	\end{frame}
	\begin{frame}{Answer 28}
		\textbf{Correct: A} --- log\_model can create the registered model version directly.\\
		\vspace{0.5em}
		\textbf{Wrong:}
		\begin{itemize}
			\item B: not a parameter.
			\item C: saves locally only.
			\item D: only logs a parameter.
		\end{itemize}
	\end{frame}
	
	% ---------------------------------------------------------------
	\begin{frame}{Question 29}
		Which call registers a model from an existing run?
		\begin{enumerate}[label=\Alph*.]
			\item \texttt{mlflow.sklearn.save\_model(run\_id)}
			\item \texttt{mlflow.create\_model(run\_id)}
			\item \texttt{mlflow.log\_metric("register", 1)}
			\item \texttt{mlflow.register\_model("runs:/<run\_id>/model", "catalog.schema.name")}
		\end{enumerate}
	\end{frame}
	\begin{frame}{Answer 29}
		\textbf{Correct: D} --- it creates a model version from the run's model artifact.\\
		\vspace{0.5em}
		\textbf{Wrong:}
		\begin{itemize}
			\item A: saves locally.
			\item B: not an API.
			\item C: logs a metric only.
		\end{itemize}
	\end{frame}
	
	% ---------------------------------------------------------------
	\begin{frame}{Question 30}
		Which URI format loads a model directly from a run?
		\begin{enumerate}[label=\Alph*.]
			\item \texttt{dbfs:/<run\_id>}
			\item \texttt{runs:/<run\_id>/model}
			\item \texttt{experiments:/<id>/model}
			\item \texttt{run://model/<run\_id>}
		\end{enumerate}
	\end{frame}
	\begin{frame}{Answer 30}
		\textbf{Correct: B} --- this is the run-based model URI.\\
		\vspace{0.5em}
		\textbf{Wrong:}
		\begin{itemize}
			\item A: points to storage, not the model URI scheme.
			\item C: not a model URI scheme.
			\item D: not valid.
		\end{itemize}
	\end{frame}
	
	% ---------------------------------------------------------------
	\begin{frame}{Question 31}
		Which call points MLflow at the Unity Catalog model registry?
		\begin{enumerate}[label=\Alph*.]
			\item \texttt{mlflow.set\_registry\_uri("uc://")}
			\item \texttt{mlflow.set\_registry\_uri("databricks")}
			\item \texttt{mlflow.set\_tracking\_uri("unity")}
			\item \texttt{mlflow.set\_registry\_uri("databricks-uc")}
		\end{enumerate}
	\end{frame}
	\begin{frame}{Answer 31}
		\textbf{Correct: D} --- this selects UC as the registry.\\
		\vspace{0.5em}
		\textbf{Wrong:}
		\begin{itemize}
			\item A: not valid.
			\item B: selects the legacy workspace registry.
			\item C: tracking URI is not the registry.
		\end{itemize}
	\end{frame}
	
	% ---------------------------------------------------------------
	\begin{frame}{Question 32}
		What replaces the legacy Staging/Production stages in the Unity Catalog model registry?
		\begin{enumerate}[label=\Alph*.]
			\item Aliases and tags
			\item Experiment folders
			\item Git branches
			\item Cluster tags
		\end{enumerate}
	\end{frame}
	\begin{frame}{Answer 32}
		\textbf{Correct: A} --- UC uses aliases such as Champion to mark versions.\\
		\vspace{0.5em}
		\textbf{Wrong:}
		\begin{itemize}
			\item B: do not label versions.
			\item C: not part of the registry.
			\item D: label compute, not models.
		\end{itemize}
	\end{frame}
	
	% ---------------------------------------------------------------
	\begin{frame}{Question 33}
		Which URI loads the model version that currently has the alias Champion?
		\begin{enumerate}[label=\Alph*.]
			\item \texttt{runs:/Champion}
			\item \texttt{models:/main.ml.churn/Champion}
			\item \texttt{models:/main.ml.churn@Champion}
			\item \texttt{models:/Champion/main.ml.churn}
		\end{enumerate}
	\end{frame}
	\begin{frame}{Answer 33}
		\textbf{Correct: C} --- the @ syntax resolves an alias.\\
		\vspace{0.5em}
		\textbf{Wrong:}
		\begin{itemize}
			\item A: not a run ID.
			\item B: legacy stage syntax.
			\item D: wrong order.
		\end{itemize}
	\end{frame}
	
	% ---------------------------------------------------------------
	\begin{frame}{Question 34}
		Which call assigns the alias Champion to version 3 of a UC model?
		\begin{enumerate}[label=\Alph*.]
			\item \texttt{client.transition\_model\_version\_stage(name, 3, "Champion")}
			\item \texttt{client.register\_alias(3)}
			\item \texttt{MlflowClient().set\_registered\_model\_alias(name, "Champion", 3)}
			\item \texttt{mlflow.set\_alias(name, 3)}
		\end{enumerate}
	\end{frame}
	\begin{frame}{Answer 34}
		\textbf{Correct: C} --- this sets an alias on a specific version.\\
		\vspace{0.5em}
		\textbf{Wrong:}
		\begin{itemize}
			\item A: stages are legacy.
			\item B: not an API.
			\item D: not an API.
		\end{itemize}
	\end{frame}
	
	% ---------------------------------------------------------------
	\begin{frame}{Question 35}
		Production loads \texttt{models:/main.ml.churn@Champion}. A new version is validated. How do you promote it with no code change?
		\begin{enumerate}[label=\Alph*.]
			\item Edit every notebook with the new version number
			\item Delete the registered model
			\item Reassign the Champion alias to the new version
			\item Rename the experiment
		\end{enumerate}
	\end{frame}
	\begin{frame}{Answer 35}
		\textbf{Correct: C} --- clients resolve the alias at load time.\\
		\vspace{0.5em}
		\textbf{Wrong:}
		\begin{itemize}
			\item A: unnecessary work.
			\item B: breaks consumers.
			\item D: has no effect on aliases.
		\end{itemize}
	\end{frame}
	
	% ---------------------------------------------------------------
	\begin{frame}{Question 36}
		Which privilege lets a group load and call a registered model in Unity Catalog (in addition to USE CATALOG and USE SCHEMA)?
		\begin{enumerate}[label=\Alph*.]
			\item \texttt{SELECT} on the model
			\item \texttt{CREATE TABLE} on the schema
			\item \texttt{EXECUTE} on the model
			\item \texttt{MODIFY} on the cluster
		\end{enumerate}
	\end{frame}
	\begin{frame}{Answer 36}
		\textbf{Correct: C} --- EXECUTE allows using the model.\\
		\vspace{0.5em}
		\textbf{Wrong:}
		\begin{itemize}
			\item A: applies to tables.
			\item B: for tables.
			\item D: unrelated.
		\end{itemize}
	\end{frame}
	
	% ---------------------------------------------------------------
	\begin{frame}{Question 37}
		Which client is used for feature tables governed by Unity Catalog?
		\begin{enumerate}[label=\Alph*.]
			\item \texttt{FeatureEngineeringClient}
			\item \texttt{MlflowClient}
			\item \texttt{FeatureStoreClient} (workspace only)
			\item \texttt{SparkSession}
		\end{enumerate}
	\end{frame}
	\begin{frame}{Answer 37}
		\textbf{Correct: A} --- from \texttt{databricks.feature\_engineering}.\\
		\vspace{0.5em}
		\textbf{Wrong:}
		\begin{itemize}
			\item B: manages tracking and registry.
			\item C: the older workspace feature store client.
			\item D: is Spark, not a feature client.
		\end{itemize}
	\end{frame}
	
	% ---------------------------------------------------------------
	\begin{frame}{Question 38}
		What does a feature table require?
		\begin{enumerate}[label=\Alph*.]
			\item An MLflow run
			\item A serving endpoint
			\item A GPU cluster
			\item A primary key
		\end{enumerate}
	\end{frame}
	\begin{frame}{Answer 38}
		\textbf{Correct: D} --- rows are identified by the primary key for lookups and merges.\\
		\vspace{0.5em}
		\textbf{Wrong:}
		\begin{itemize}
			\item A: not required.
			\item B: not required.
			\item C: not required.
		\end{itemize}
	\end{frame}
	
	% ---------------------------------------------------------------
	\begin{frame}{Question 39}
		Which call creates a feature table from a DataFrame?
		\begin{enumerate}[label=\Alph*.]
			\item \texttt{fe.create\_table(name=..., primary\_keys=["customer\_id"], df=df)}
			\item \texttt{fe.new\_features(df)}
			\item \texttt{fe.log\_model(df)}
			\item \texttt{fe.score\_batch(df)}
		\end{enumerate}
	\end{frame}
	\begin{frame}{Answer 39}
		\textbf{Correct: A} --- it defines the key and initial data.\\
		\vspace{0.5em}
		\textbf{Wrong:}
		\begin{itemize}
			\item B: not an API.
			\item C: logs a model.
			\item D: scores data.
		\end{itemize}
	\end{frame}
	
	% ---------------------------------------------------------------
	\begin{frame}{Question 40}
		Why add a timestamp key to a feature table?
		\begin{enumerate}[label=\Alph*.]
			\item To enable point-in-time lookups and avoid leakage
			\item To encrypt the table
			\item To replace the primary key
			\item To make training faster
		\end{enumerate}
	\end{frame}
	\begin{frame}{Answer 40}
		\textbf{Correct: A} --- training rows join to the latest feature values as of their timestamp.\\
		\vspace{0.5em}
		\textbf{Wrong:}
		\begin{itemize}
			\item B: not its purpose.
			\item C: it is used together with it.
			\item D: not its purpose.
		\end{itemize}
	\end{frame}
	
	% ---------------------------------------------------------------
	\begin{frame}{Question 41}
		What is \texttt{FeatureLookup} used for?
		\begin{enumerate}[label=\Alph*.]
			\item Specifying which features to join to a label DataFrame by key
			\item Starting an AutoML run
			\item Registering a model
			\item Creating a cluster
		\end{enumerate}
	\end{frame}
	\begin{frame}{Answer 41}
		\textbf{Correct: A} --- it names the table, lookup key, and features.\\
		\vspace{0.5em}
		\textbf{Wrong:}
		\begin{itemize}
			\item B: not its role.
			\item C: not its role.
			\item D: not its role.
		\end{itemize}
	\end{frame}
	
	% ---------------------------------------------------------------
	\begin{frame}{Question 42}
		Which call joins feature lookups to a label DataFrame to build training data?
		\begin{enumerate}[label=\Alph*.]
			\item \texttt{mlflow.search\_runs(df)}
			\item \texttt{fe.score\_batch(df)}
			\item \texttt{fe.create\_training\_set(df, feature\_lookups, label)} and then \texttt{.load\_df()}
			\item \texttt{fe.write\_table(df)}
		\end{enumerate}
	\end{frame}
	\begin{frame}{Answer 42}
		\textbf{Correct: C} --- it assembles features and labels for training.\\
		\vspace{0.5em}
		\textbf{Wrong:}
		\begin{itemize}
			\item A: searches runs.
			\item B: performs inference.
			\item D: writes feature data.
		\end{itemize}
	\end{frame}
	
	% ---------------------------------------------------------------
	\begin{frame}{Question 43}
		Why log a model with \texttt{fe.log\_model()} instead of a plain MLflow flavor?
		\begin{enumerate}[label=\Alph*.]
			\item It disables autologging
			\item It skips the registry
			\item It stores feature lookup metadata so features are retrieved automatically at inference
			\item It makes the model smaller
		\end{enumerate}
	\end{frame}
	\begin{frame}{Answer 43}
		\textbf{Correct: C} --- this prevents training-serving skew.\\
		\vspace{0.5em}
		\textbf{Wrong:}
		\begin{itemize}
			\item A: not its purpose.
			\item B: not its purpose.
			\item D: not its purpose.
		\end{itemize}
	\end{frame}
	
	% ---------------------------------------------------------------
	\begin{frame}{Question 44}
		A model was logged with \texttt{fe.log\_model()}. At scoring time the input has only \texttt{customer\_id}. Which call should be used?
		\begin{enumerate}[label=\Alph*.]
			\item \texttt{mlflow.search\_runs(df)}
			\item \texttt{fe.create\_table(df)}
			\item \texttt{fe.score\_batch(model\_uri=..., df=df)}
			\item \texttt{model.predict(df)} with the raw key
		\end{enumerate}
	\end{frame}
	\begin{frame}{Answer 44}
		\textbf{Correct: C} --- it looks up the needed features automatically.\\
		\vspace{0.5em}
		\textbf{Wrong:}
		\begin{itemize}
			\item A: not scoring.
			\item B: creates a table.
			\item D: missing features would cause an error.
		\end{itemize}
	\end{frame}
	
	% ---------------------------------------------------------------
	\begin{frame}{Question 45}
		Which write mode updates existing feature rows by primary key and inserts new ones?
		\begin{enumerate}[label=\Alph*.]
			\item \texttt{mode="merge"}
			\item \texttt{mode="append"}
			\item \texttt{mode="delete"}
			\item \texttt{mode="overwrite"}
		\end{enumerate}
	\end{frame}
	\begin{frame}{Answer 45}
		\textbf{Correct: A} --- merge performs an upsert on the key.\\
		\vspace{0.5em}
		\textbf{Wrong:}
		\begin{itemize}
			\item B: not the supported behavior here.
			\item C: not a write mode.
			\item D: replaces the whole table.
		\end{itemize}
	\end{frame}
	
	% ---------------------------------------------------------------
	\begin{frame}{Question 46}
		What is a core benefit of using a feature store?
		\begin{enumerate}[label=\Alph*.]
			\item Automatic hyperparameter tuning
			\item Free GPU time
			\item Removing the need for labels
			\item Reusing consistent features across training and inference with lineage
		\end{enumerate}
	\end{frame}
	\begin{frame}{Answer 46}
		\textbf{Correct: D} --- this reduces duplicate work and skew.\\
		\vspace{0.5em}
		\textbf{Wrong:}
		\begin{itemize}
			\item A: that is a tuning tool.
			\item B: not a benefit.
			\item C: labels are still required.
		\end{itemize}
	\end{frame}
	
	% ---------------------------------------------------------------
	\begin{frame}{Question 47}
		A model needs low-latency feature lookups at real-time inference. What is typically required?
		\begin{enumerate}[label=\Alph*.]
			\item Publish the feature table to an online store (online table)
			\item Read the offline Delta table on every request
			\item Disable the primary key
			\item Store features in a notebook variable
		\end{enumerate}
	\end{frame}
	\begin{frame}{Answer 47}
		\textbf{Correct: A} --- online stores serve low-latency lookups.\\
		\vspace{0.5em}
		\textbf{Wrong:}
		\begin{itemize}
			\item B: too slow.
			\item C: keys are needed.
			\item D: not durable.
		\end{itemize}
	\end{frame}
	
	% ---------------------------------------------------------------
	\begin{frame}{Question 48}
		Which Databricks utility produces summary statistics and distributions for a DataFrame?
		\begin{enumerate}[label=\Alph*.]
			\item \texttt{dbutils.notebook.run(df)}
			\item \texttt{dbutils.data.summarize(df)}
			\item \texttt{dbutils.widgets.get(df)}
			\item \texttt{dbutils.fs.ls(df)}
		\end{enumerate}
	\end{frame}
	\begin{frame}{Answer 48}
		\textbf{Correct: B} --- it renders a data profile.\\
		\vspace{0.5em}
		\textbf{Wrong:}
		\begin{itemize}
			\item A: runs notebooks.
			\item C: reads widgets.
			\item D: lists files.
		\end{itemize}
	\end{frame}
	
	% ---------------------------------------------------------------
	\begin{frame}{Question 49}
		Which DataFrame method returns count, mean, stddev, min, quartiles, and max?
		\begin{enumerate}[label=\Alph*.]
			\item \texttt{df.summary()}
			\item \texttt{df.describe()}
			\item \texttt{df.explain()}
			\item \texttt{df.printSchema()}
		\end{enumerate}
	\end{frame}
	\begin{frame}{Answer 49}
		\textbf{Correct: A} --- summary includes percentiles, unlike describe.\\
		\vspace{0.5em}
		\textbf{Wrong:}
		\begin{itemize}
			\item B: omits the quartiles by default.
			\item C: shows the query plan.
			\item D: shows column types.
		\end{itemize}
	\end{frame}
	
	% ---------------------------------------------------------------
	\begin{frame}{Question 50}
		Which statistic is most robust for imputing a skewed numeric feature?
		\begin{enumerate}[label=\Alph*.]
			\item Median
			\item Maximum
			\item Mean
			\item Sum
		\end{enumerate}
	\end{frame}
	\begin{frame}{Answer 50}
		\textbf{Correct: A} --- the median is not pulled by extreme values.\\
		\vspace{0.5em}
		\textbf{Wrong:}
		\begin{itemize}
			\item B: distorts the distribution.
			\item C: is pulled by outliers.
			\item D: is not a valid imputation value.
		\end{itemize}
	\end{frame}
	
	% ---------------------------------------------------------------
	\begin{frame}{Question 51}
		Which strategies does \texttt{pyspark.ml.feature.Imputer} support?
		\begin{enumerate}[label=\Alph*.]
			\item forward-fill and back-fill
			\item regression and KNN
			\item mean, max, and min
			\item mean, median, and mode
		\end{enumerate}
	\end{frame}
	\begin{frame}{Answer 51}
		\textbf{Correct: D} --- these are the available strategies.\\
		\vspace{0.5em}
		\textbf{Wrong:}
		\begin{itemize}
			\item A: not supported.
			\item B: not supported.
			\item C: max and min are not strategies.
		\end{itemize}
	\end{frame}
	
	% ---------------------------------------------------------------
	\begin{frame}{Question 52}
		A team wants to preserve information about which values were missing. What should they add?
		\begin{enumerate}[label=\Alph*.]
			\item Drop the column
			\item Impute and discard the original
			\item Cast all values to string
			\item A boolean indicator column for missingness before imputing
		\end{enumerate}
	\end{frame}
	\begin{frame}{Answer 52}
		\textbf{Correct: D} --- the model can learn from the missingness pattern.\\
		\vspace{0.5em}
		\textbf{Wrong:}
		\begin{itemize}
			\item A: loses information.
			\item B: loses missingness information.
			\item C: does not help.
		\end{itemize}
	\end{frame}
	
	% ---------------------------------------------------------------
	\begin{frame}{Question 53}
		A data scientist calls \texttt{transform()} directly on a new \texttt{Imputer} and gets an error. Why?
		\begin{enumerate}[label=\Alph*.]
			\item \texttt{transform()} needs a GPU
			\item \texttt{Imputer} is an estimator and must be fit first
			\item \texttt{Imputer} requires a feature table
			\item \texttt{Imputer} only works on strings
		\end{enumerate}
	\end{frame}
	\begin{frame}{Answer 53}
		\textbf{Correct: B} --- \texttt{fit()} returns an \texttt{ImputerModel}, which has \texttt{transform()}.\\
		\vspace{0.5em}
		\textbf{Wrong:}
		\begin{itemize}
			\item A: not true.
			\item C: not true.
			\item D: it works on numeric columns.
		\end{itemize}
	\end{frame}
	
	% ---------------------------------------------------------------
	\begin{frame}{Question 54}
		Which is the correct Spark ML order to one-hot encode a string column?
		\begin{enumerate}[label=\Alph*.]
			\item \texttt{VectorAssembler} then \texttt{StringIndexer}
			\item \texttt{Imputer} then \texttt{OneHotEncoder} only
			\item \texttt{OneHotEncoder} then \texttt{StringIndexer}
			\item \texttt{StringIndexer} then \texttt{OneHotEncoder}
		\end{enumerate}
	\end{frame}
	\begin{frame}{Answer 54}
		\textbf{Correct: D} --- strings must be indexed to numbers before encoding.\\
		\vspace{0.5em}
		\textbf{Wrong:}
		\begin{itemize}
			\item A: assembler builds the final vector.
			\item B: strings are not indexed.
			\item C: encoder needs numeric indices.
		\end{itemize}
	\end{frame}
	
	% ---------------------------------------------------------------
	\begin{frame}{Question 55}
		What does \texttt{StringIndexer(handleInvalid="keep")} do for labels unseen during training?
		\begin{enumerate}[label=\Alph*.]
			\item Raises an error
			\item Deletes the column
			\item Maps them to the mean
			\item Assigns them an extra index instead of failing
		\end{enumerate}
	\end{frame}
	\begin{frame}{Answer 55}
		\textbf{Correct: D} --- this avoids errors on new categories at inference.\\
		\vspace{0.5em}
		\textbf{Wrong:}
		\begin{itemize}
			\item A: that is the default \texttt{error}.
			\item B: not the behavior.
			\item C: not the behavior.
		\end{itemize}
	\end{frame}
	
	% ---------------------------------------------------------------
	\begin{frame}{Question 56}
		What does \texttt{VectorAssembler} do?
		\begin{enumerate}[label=\Alph*.]
			\item Scales features to unit variance
			\item Trains a model
			\item Splits a vector into columns
			\item Combines selected columns into a single vector column
		\end{enumerate}
	\end{frame}
	\begin{frame}{Answer 56}
		\textbf{Correct: D} --- Spark ML models expect one features vector.\\
		\vspace{0.5em}
		\textbf{Wrong:}
		\begin{itemize}
			\item A: that is \texttt{StandardScaler}.
			\item B: it is a transformer.
			\item C: the opposite.
		\end{itemize}
	\end{frame}
	
	% ---------------------------------------------------------------
	\begin{frame}{Question 57}
		What does \texttt{Pipeline.fit()} return?
		\begin{enumerate}[label=\Alph*.]
			\item A pandas DataFrame
			\item A \texttt{Pipeline}
			\item A \texttt{CrossValidator}
			\item A \texttt{PipelineModel}
		\end{enumerate}
	\end{frame}
	\begin{frame}{Answer 57}
		\textbf{Correct: D} --- fitting estimators produces fitted transformers chained as one model.\\
		\vspace{0.5em}
		\textbf{Wrong:}
		\begin{itemize}
			\item A: not what it returns.
			\item B: that is the unfitted object.
			\item C: not what it returns.
		\end{itemize}
	\end{frame}
	
	% ---------------------------------------------------------------
	\begin{frame}{Question 58}
		Which model type generally needs feature scaling the least?
		\begin{enumerate}[label=\Alph*.]
			\item Tree-based models such as random forests
			\item Neural networks
			\item Linear regression with regularization
			\item k-nearest neighbors
		\end{enumerate}
	\end{frame}
	\begin{frame}{Answer 58}
		\textbf{Correct: A} --- splits depend on order, not magnitude.\\
		\vspace{0.5em}
		\textbf{Wrong:}
		\begin{itemize}
			\item B: typically benefit from scaling.
			\item C: regularization is scale sensitive.
			\item D: uses distances and benefits from scaling.
		\end{itemize}
	\end{frame}
	
	% ---------------------------------------------------------------
	\begin{frame}{Question 59}
		A model is trained on log(price). What must be done before computing RMSE against original prices?
		\begin{enumerate}[label=\Alph*.]
			\item Back-transform predictions with exp
			\item Nothing, compare directly
			\item Drop the highest predictions
			\item Take the log of the original prices twice
		\end{enumerate}
	\end{frame}
	\begin{frame}{Answer 59}
		\textbf{Correct: A} --- metrics must compare values on the same scale.\\
		\vspace{0.5em}
		\textbf{Wrong:}
		\begin{itemize}
			\item B: the scales differ.
			\item C: not a fix.
			\item D: wrong transformation.
		\end{itemize}
	\end{frame}
	
	% ---------------------------------------------------------------
	\begin{frame}{Question 60}
		What is lazy evaluation in Spark?
		\begin{enumerate}[label=\Alph*.]
			\item All operations run immediately
			\item Only SQL is optimized
			\item Transformations are not run until an action is called
			\item Data is deleted after each step
		\end{enumerate}
	\end{frame}
	\begin{frame}{Answer 60}
		\textbf{Correct: C} --- actions such as \texttt{count()} trigger execution.\\
		\vspace{0.5em}
		\textbf{Wrong:}
		\begin{itemize}
			\item A: that is eager.
			\item B: not true.
			\item D: not true.
		\end{itemize}
	\end{frame}
	
	% ---------------------------------------------------------------
	\begin{frame}{Question 61}
		Which call converts a Spark DataFrame to a pandas API on Spark DataFrame, keeping it distributed?
		\begin{enumerate}[label=\Alph*.]
			\item \texttt{df.pandas\_api()}
			\item \texttt{df.collect()}
			\item \texttt{df.toPandas()}
			\item \texttt{df.cache()}
		\end{enumerate}
	\end{frame}
	\begin{frame}{Answer 61}
		\textbf{Correct: A} --- it returns a distributed pandas-style DataFrame.\\
		\vspace{0.5em}
		\textbf{Wrong:}
		\begin{itemize}
			\item B: returns a list on the driver.
			\item C: collects everything to the driver.
			\item D: only caches.
		\end{itemize}
	\end{frame}
	
	% ---------------------------------------------------------------
	\begin{frame}{Question 62}
		Why is \texttt{toPandas()} risky on a very large DataFrame?
		\begin{enumerate}[label=\Alph*.]
			\item It converts to RDD only
			\item It collects all data to the driver and may run out of memory
			\item It deletes the Spark table
			\item It always runs faster
		\end{enumerate}
	\end{frame}
	\begin{frame}{Answer 62}
		\textbf{Correct: B} --- the driver holds the entire result.\\
		\vspace{0.5em}
		\textbf{Wrong:}
		\begin{itemize}
			\item A: not true.
			\item C: not true.
			\item D: not true.
		\end{itemize}
	\end{frame}
	
	% ---------------------------------------------------------------
	\begin{frame}{Question 63}
		Why are pandas UDFs generally faster than row-at-a-time Python UDFs?
		\begin{enumerate}[label=\Alph*.]
			\item They skip serialization
			\item They use GPUs by default
			\item They process batches using Apache Arrow
			\item They run only on the driver
		\end{enumerate}
	\end{frame}
	\begin{frame}{Answer 63}
		\textbf{Correct: C} --- vectorized execution avoids per-row overhead.\\
		\vspace{0.5em}
		\textbf{Wrong:}
		\begin{itemize}
			\item A: Arrow still transfers data.
			\item B: not by default.
			\item D: they run on workers.
		\end{itemize}
	\end{frame}
	
	% ---------------------------------------------------------------
	\begin{frame}{Question 64}
		A team wants to train one model per store in parallel across a cluster. What should they use?
		\begin{enumerate}[label=\Alph*.]
			\item \texttt{df.coalesce(1)}
			\item \texttt{df.groupBy("store").applyInPandas(train\_fn, schema)}
			\item \texttt{df.toPandas()} and a for-loop
			\item \texttt{df.limit(1)}
		\end{enumerate}
	\end{frame}
	\begin{frame}{Answer 64}
		\textbf{Correct: B} --- each group is processed as a pandas DataFrame on workers.\\
		\vspace{0.5em}
		\textbf{Wrong:}
		\begin{itemize}
			\item A: forces one partition.
			\item C: runs on the driver.
			\item D: keeps one row.
		\end{itemize}
	\end{frame}
	
	% ---------------------------------------------------------------
	\begin{frame}{Question 65}
		Which call performs a reproducible random train/test split of a Spark DataFrame?
		\begin{enumerate}[label=\Alph*.]
			\item \texttt{train\_test\_split(df)} from scikit-learn on a Spark DataFrame
			\item \texttt{df.split(0.8)}
			\item \texttt{df.limit(0.8)}
			\item \texttt{df.randomSplit([0.8, 0.2], seed=42)}
		\end{enumerate}
	\end{frame}
	\begin{frame}{Answer 65}
		\textbf{Correct: D} --- the seed makes the split repeatable.\\
		\vspace{0.5em}
		\textbf{Wrong:}
		\begin{itemize}
			\item A: expects pandas or arrays.
			\item B: not an API.
			\item C: limit takes a row count.
		\end{itemize}
	\end{frame}
	
	% ---------------------------------------------------------------
	\begin{frame}{Question 66}
		Which step is an example of data leakage?
		\begin{enumerate}[label=\Alph*.]
			\item Using a fixed random seed
			\item Computing normalization statistics on the full dataset before splitting
			\item Splitting before preprocessing
			\item Fitting the scaler on training data only
		\end{enumerate}
	\end{frame}
	\begin{frame}{Answer 66}
		\textbf{Correct: B} --- test information leaks into training.\\
		\vspace{0.5em}
		\textbf{Wrong:}
		\begin{itemize}
			\item A: this aids reproducibility.
			\item C: this is correct practice.
			\item D: this is correct practice.
		\end{itemize}
	\end{frame}
	
	% ---------------------------------------------------------------
	\begin{frame}{Question 67}
		A dataset has 1 percent positive labels. Which split technique keeps class proportions?
		\begin{enumerate}[label=\Alph*.]
			\item Sorting by label
			\item Stratified splitting
			\item Using only the first rows
			\item Splitting by column order
		\end{enumerate}
	\end{frame}
	\begin{frame}{Answer 67}
		\textbf{Correct: B} --- stratification preserves the label ratio.\\
		\vspace{0.5em}
		\textbf{Wrong:}
		\begin{itemize}
			\item A: creates biased splits.
			\item C: is not representative.
			\item D: ignores class balance.
		\end{itemize}
	\end{frame}
	
	% ---------------------------------------------------------------
	\begin{frame}{Question 68}
		How should a time series dataset be split for evaluation?
		\begin{enumerate}[label=\Alph*.]
			\item Train on earlier data and test on later data without shuffling
			\item Random shuffle split
			\item Use the same rows for train and test
			\item Train on the latest and test on the earliest
		\end{enumerate}
	\end{frame}
	\begin{frame}{Answer 68}
		\textbf{Correct: A} --- this mimics real forecasting and avoids future leakage.\\
		\vspace{0.5em}
		\textbf{Wrong:}
		\begin{itemize}
			\item B: leaks future information.
			\item C: no unseen test data.
			\item D: reverses time.
		\end{itemize}
	\end{frame}
	
	% ---------------------------------------------------------------
	\begin{frame}{Question 69}
		Which technique helps with a heavily imbalanced classification dataset?
		\begin{enumerate}[label=\Alph*.]
			\item Using accuracy as the only metric
			\item Class weights or resampling the training data
			\item Duplicating the test set
			\item Removing the minority class
		\end{enumerate}
	\end{frame}
	\begin{frame}{Answer 69}
		\textbf{Correct: B} --- these reduce bias toward the majority class.\\
		\vspace{0.5em}
		\textbf{Wrong:}
		\begin{itemize}
			\item A: accuracy is misleading.
			\item C: contaminates evaluation.
			\item D: destroys the signal.
		\end{itemize}
	\end{frame}
	
	% ---------------------------------------------------------------
	\begin{frame}{Question 70}
		Hyperopt \texttt{fmin} minimizes the objective. How do you maximize accuracy?
		\begin{enumerate}[label=\Alph*.]
			\item Return accuracy directly
			\item Return the negative accuracy as the loss
			\item Use \texttt{algo=max.suggest}
			\item Set \texttt{maximize=True}
		\end{enumerate}
	\end{frame}
	\begin{frame}{Answer 70}
		\textbf{Correct: B} --- minimizing -accuracy maximizes accuracy.\\
		\vspace{0.5em}
		\textbf{Wrong:}
		\begin{itemize}
			\item A: this would minimize accuracy.
			\item C: not an algorithm.
			\item D: no such argument.
		\end{itemize}
	\end{frame}
	
	% ---------------------------------------------------------------
	\begin{frame}{Question 71}
		When is \texttt{SparkTrials} an appropriate choice with Hyperopt?
		\begin{enumerate}[label=\Alph*.]
			\item When using no objective function
			\item When tuning only on the driver
			\item When tuning single-machine libraries such as scikit-learn
			\item When tuning distributed Spark ML models
		\end{enumerate}
	\end{frame}
	\begin{frame}{Answer 71}
		\textbf{Correct: C} --- each trial runs on a worker in parallel.\\
		\vspace{0.5em}
		\textbf{Wrong:}
		\begin{itemize}
			\item A: an objective is required.
			\item B: SparkTrials uses workers.
			\item D: those already use the cluster.
		\end{itemize}
	\end{frame}
	
	% ---------------------------------------------------------------
	\begin{frame}{Question 72}
		What does the \texttt{parallelism} setting in \texttt{SparkTrials} control?
		\begin{enumerate}[label=\Alph*.]
			\item The number of hyperparameters
			\item The model's tree depth
			\item The maximum number of trials evaluated concurrently
			\item The number of folds
		\end{enumerate}
	\end{frame}
	\begin{frame}{Answer 72}
		\textbf{Correct: C} --- more parallelism is faster but makes TPE less adaptive.\\
		\vspace{0.5em}
		\textbf{Wrong:}
		\begin{itemize}
			\item A: not controlled here.
			\item B: not controlled here.
			\item D: that is for cross-validation.
		\end{itemize}
	\end{frame}
	
	% ---------------------------------------------------------------
	\begin{frame}{Question 73}
		Which Hyperopt space is best for an integer-like parameter such as \texttt{max\_depth}?
		\begin{enumerate}[label=\Alph*.]
			\item \texttt{hp.uniform("max\_depth", 2, 20)}
			\item \texttt{hp.quniform("max\_depth", 2, 20, 1)} with an int cast
			\item \texttt{hp.pyll("max\_depth")}
			\item \texttt{hp.normal("max\_depth", 0, 1)}
		\end{enumerate}
	\end{frame}
	\begin{frame}{Answer 73}
		\textbf{Correct: B} --- quniform returns rounded floats, so cast to int.\\
		\vspace{0.5em}
		\textbf{Wrong:}
		\begin{itemize}
			\item A: returns arbitrary floats.
			\item C: not an API.
			\item D: wrong distribution.
		\end{itemize}
	\end{frame}
	
	% ---------------------------------------------------------------
	\begin{frame}{Question 74}
		\texttt{fmin} is run with \texttt{hp.choice} for a parameter. What does the returned best dict contain for it?
		\begin{enumerate}[label=\Alph*.]
			\item A Spark DataFrame
			\item The loss value
			\item The actual chosen value
			\item The index of the chosen option
		\end{enumerate}
	\end{frame}
	\begin{frame}{Answer 74}
		\textbf{Correct: D} --- use \texttt{space\_eval} to recover the actual value.\\
		\vspace{0.5em}
		\textbf{Wrong:}
		\begin{itemize}
			\item A: not what it returns.
			\item B: not for the parameter.
			\item C: not for hp.choice.
		\end{itemize}
	\end{frame}
	
	% ---------------------------------------------------------------
	\begin{frame}{Question 75}
		How does \texttt{tpe.suggest} differ from \texttt{rand.suggest}?
		\begin{enumerate}[label=\Alph*.]
			\item TPE only works for Spark ML
			\item TPE samples uniformly at random
			\item TPE uses past trial results to choose promising values (Bayesian-style)
			\item TPE does not need an objective
		\end{enumerate}
	\end{frame}
	\begin{frame}{Answer 75}
		\textbf{Correct: C} --- random search ignores previous results.\\
		\vspace{0.5em}
		\textbf{Wrong:}
		\begin{itemize}
			\item A: it works for any objective.
			\item B: that is rand.suggest.
			\item D: it needs one.
		\end{itemize}
	\end{frame}
	
	% ---------------------------------------------------------------
	\begin{frame}{Question 76}
		A 3-fold CrossValidator uses a grid of 4 x 3 = 12 parameter combinations. How many models are trained during the search (excluding the final refit)?
		\begin{enumerate}[label=\Alph*.]
			\item 36
			\item 12
			\item 15
			\item 3
		\end{enumerate}
	\end{frame}
	\begin{frame}{Answer 76}
		\textbf{Correct: A} --- 12 combinations times 3 folds.\\
		\vspace{0.5em}
		\textbf{Wrong:}
		\begin{itemize}
			\item B: ignores the folds.
			\item C: adds instead of multiplying.
			\item D: counts folds only.
		\end{itemize}
	\end{frame}
	
	% ---------------------------------------------------------------
	\begin{frame}{Question 77}
		How does \texttt{TrainValidationSplit} differ from \texttt{CrossValidator}?
		\begin{enumerate}[label=\Alph*.]
			\item It uses k folds
			\item It only works for clustering
			\item It evaluates each parameter combination once on a single split
			\item It tunes without an evaluator
		\end{enumerate}
	\end{frame}
	\begin{frame}{Answer 77}
		\textbf{Correct: C} --- it is cheaper but less reliable.\\
		\vspace{0.5em}
		\textbf{Wrong:}
		\begin{itemize}
			\item A: that is \texttt{CrossValidator}.
			\item B: not true.
			\item D: an evaluator is needed.
		\end{itemize}
	\end{frame}
	
	% ---------------------------------------------------------------
	\begin{frame}{Question 78}
		Which three things must be supplied when constructing a Spark ML \texttt{CrossValidator}?
		\begin{enumerate}[label=\Alph*.]
			\item feature table, lookup, and label
			\item model, registry, and endpoint
			\item cluster, runtime, and job
			\item estimator, estimatorParamMaps, and evaluator
		\end{enumerate}
	\end{frame}
	\begin{frame}{Answer 78}
		\textbf{Correct: D} --- the estimator, grid, and metric define the search.\\
		\vspace{0.5em}
		\textbf{Wrong:}
		\begin{itemize}
			\item A: unrelated.
			\item B: unrelated.
			\item C: unrelated.
		\end{itemize}
	\end{frame}
	
	% ---------------------------------------------------------------
	\begin{frame}{Question 79}
		A Spark ML decision tree fails because a categorical feature has more categories than allowed. What should be changed?
		\begin{enumerate}[label=\Alph*.]
			\item Increase \texttt{numTrees} only
			\item Increase \texttt{maxBins}
			\item Disable the pipeline
			\item Decrease \texttt{maxDepth}
		\end{enumerate}
	\end{frame}
	\begin{frame}{Answer 79}
		\textbf{Correct: B} --- maxBins must be at least the number of categories.\\
		\vspace{0.5em}
		\textbf{Wrong:}
		\begin{itemize}
			\item A: does not address bins.
			\item C: not a fix.
			\item D: does not address bins.
		\end{itemize}
	\end{frame}
	
	% ---------------------------------------------------------------
	\begin{frame}{Question 80}
		A fraud dataset has 1 percent positives. Why is accuracy a poor metric?
		\begin{enumerate}[label=\Alph*.]
			\item Accuracy ignores true negatives
			\item Predicting all negatives gives 99 percent accuracy
			\item Accuracy cannot be computed
			\item Accuracy is only for regression
		\end{enumerate}
	\end{frame}
	\begin{frame}{Answer 80}
		\textbf{Correct: B} --- accuracy hides failure on the rare class.\\
		\vspace{0.5em}
		\textbf{Wrong:}
		\begin{itemize}
			\item A: it includes them.
			\item C: it can be computed.
			\item D: it is for classification.
		\end{itemize}
	\end{frame}
	
	% ---------------------------------------------------------------
	\begin{frame}{Question 81}
		A screening model must catch as many true positives as possible. Which metric should be maximized?
		\begin{enumerate}[label=\Alph*.]
			\item Recall
			\item Precision
			\item RMSE
			\item Silhouette score
		\end{enumerate}
	\end{frame}
	\begin{frame}{Answer 81}
		\textbf{Correct: A} --- recall measures the share of positives found.\\
		\vspace{0.5em}
		\textbf{Wrong:}
		\begin{itemize}
			\item B: measures correctness of positive predictions.
			\item C: a regression metric.
			\item D: a clustering metric.
		\end{itemize}
	\end{frame}
	
	% ---------------------------------------------------------------
	\begin{frame}{Question 82}
		What does ROC AUC represent?
		\begin{enumerate}[label=\Alph*.]
			\item The fraction of correct predictions
			\item The number of features used
			\item The average squared error
			\item The probability that a random positive is ranked above a random negative
		\end{enumerate}
	\end{frame}
	\begin{frame}{Answer 82}
		\textbf{Correct: D} --- it measures ranking quality across thresholds.\\
		\vspace{0.5em}
		\textbf{Wrong:}
		\begin{itemize}
			\item A: that is accuracy.
			\item B: not related.
			\item C: that is MSE.
		\end{itemize}
	\end{frame}
	
	% ---------------------------------------------------------------
	\begin{frame}{Question 83}
		RMSE is much larger than MAE for a regression model. What does this suggest?
		\begin{enumerate}[label=\Alph*.]
			\item The model is perfect
			\item A few large errors are present
			\item There are no outliers
			\item The features are scaled
		\end{enumerate}
	\end{frame}
	\begin{frame}{Answer 83}
		\textbf{Correct: B} --- RMSE squares errors and penalizes large ones more.\\
		\vspace{0.5em}
		\textbf{Wrong:}
		\begin{itemize}
			\item A: errors are present.
			\item C: outliers are likely.
			\item D: unrelated.
		\end{itemize}
	\end{frame}
	
	% ---------------------------------------------------------------
	\begin{frame}{Question 84}
		A regression model has R2 close to 0 on the test set. What does this mean?
		\begin{enumerate}[label=\Alph*.]
			\item It has zero error
			\item It explains all variance
			\item It performs about as well as predicting the mean
			\item It is perfectly calibrated
		\end{enumerate}
	\end{frame}
	\begin{frame}{Answer 84}
		\textbf{Correct: C} --- R2 compares the model to a mean baseline.\\
		\vspace{0.5em}
		\textbf{Wrong:}
		\begin{itemize}
			\item A: not implied.
			\item B: that is R2 near 1.
			\item D: not implied.
		\end{itemize}
	\end{frame}
	
	% ---------------------------------------------------------------
	\begin{frame}{Question 85}
		Training accuracy is 99 percent and validation accuracy is 72 percent. What is the likely issue and a fix?
		\begin{enumerate}[label=\Alph*.]
			\item Label leakage in serving
			\item Underfitting; add more layers
			\item Overfitting; add regularization, simplify the model, or get more data
			\item Data drift; redeploy
		\end{enumerate}
	\end{frame}
	\begin{frame}{Answer 85}
		\textbf{Correct: C} --- a large train/validation gap indicates overfitting.\\
		\vspace{0.5em}
		\textbf{Wrong:}
		\begin{itemize}
			\item A: not indicated.
			\item B: underfitting shows poor training performance.
			\item D: not indicated.
		\end{itemize}
	\end{frame}
	
	% ---------------------------------------------------------------
	\begin{frame}{Question 86}
		Which statement correctly contrasts bagging and boosting?
		\begin{enumerate}[label=\Alph*.]
			\item Boosting trains independent models in parallel
			\item Bagging mainly reduces variance; boosting builds models sequentially to reduce bias
			\item Bagging builds models sequentially
			\item Both only reduce bias
		\end{enumerate}
	\end{frame}
	\begin{frame}{Answer 86}
		\textbf{Correct: B} --- random forests bag; gradient boosting is sequential.\\
		\vspace{0.5em}
		\textbf{Wrong:}
		\begin{itemize}
			\item A: that is bagging.
			\item C: that is boosting.
			\item D: bagging mainly reduces variance.
		\end{itemize}
	\end{frame}
	
	% ---------------------------------------------------------------
	\begin{frame}{Question 87}
		Which regularization can shrink some coefficients exactly to zero?
		\begin{enumerate}[label=\Alph*.]
			\item Dropout on linear models
			\item No regularization
			\item L1 (Lasso)
			\item L2 (Ridge)
		\end{enumerate}
	\end{frame}
	\begin{frame}{Answer 87}
		\textbf{Correct: C} --- L1 produces sparse solutions and acts as feature selection.\\
		\vspace{0.5em}
		\textbf{Wrong:}
		\begin{itemize}
			\item A: not the standard answer.
			\item B: no shrinkage.
			\item D: shrinks but rarely zeroes coefficients.
		\end{itemize}
	\end{frame}
	
	% ---------------------------------------------------------------
	\begin{frame}{Question 88}
		What is the default metric of Spark ML \texttt{RegressionEvaluator}?
		\begin{enumerate}[label=\Alph*.]
			\item areaUnderROC
			\item r2
			\item f1
			\item rmse
		\end{enumerate}
	\end{frame}
	\begin{frame}{Answer 88}
		\textbf{Correct: D} --- rmse is the default metric name.\\
		\vspace{0.5em}
		\textbf{Wrong:}
		\begin{itemize}
			\item A: belongs to the binary classification evaluator.
			\item B: available but not the default.
			\item C: belongs to the multiclass evaluator.
		\end{itemize}
	\end{frame}
	
	% ---------------------------------------------------------------
	\begin{frame}{Question 89}
		Which Spark ML evaluator computes areaUnderROC?
		\begin{enumerate}[label=\Alph*.]
			\item \texttt{BinaryClassificationEvaluator}
			\item \texttt{RankingEvaluator}
			\item \texttt{ClusteringEvaluator}
			\item \texttt{RegressionEvaluator}
		\end{enumerate}
	\end{frame}
	\begin{frame}{Answer 89}
		\textbf{Correct: A} --- its default metric is areaUnderROC.\\
		\vspace{0.5em}
		\textbf{Wrong:}
		\begin{itemize}
			\item B: not the answer here.
			\item C: clustering metrics.
			\item D: regression metrics.
		\end{itemize}
	\end{frame}
	
	% ---------------------------------------------------------------
	\begin{frame}{Question 90}
		What is the difference between a Spark ML estimator and a transformer?
		\begin{enumerate}[label=\Alph*.]
			\item Estimators transform data directly
			\item They are the same class
			\item An estimator has \texttt{fit()} and returns a model; a transformer has \texttt{transform()}
			\item Transformers have \texttt{fit()} only
		\end{enumerate}
	\end{frame}
	\begin{frame}{Answer 90}
		\textbf{Correct: C} --- fit produces a fitted transformer.\\
		\vspace{0.5em}
		\textbf{Wrong:}
		\begin{itemize}
			\item A: they must be fit first.
			\item B: they are distinct.
			\item D: they have \texttt{transform()}.
		\end{itemize}
	\end{frame}
	
	% ---------------------------------------------------------------
	\begin{frame}{Question 91}
		How can you inspect which features a fitted random forest relies on most?
		\begin{enumerate}[label=\Alph*.]
			\item Look at the learning rate
			\item Look at the cluster size
			\item Look at the run ID
			\item Look at feature importances
		\end{enumerate}
	\end{frame}
	\begin{frame}{Answer 91}
		\textbf{Correct: D} --- tree ensembles expose importances (\texttt{featureImportances} in Spark ML).\\
		\vspace{0.5em}
		\textbf{Wrong:}
		\begin{itemize}
			\item A: not importances.
			\item B: not model-related.
			\item C: not importances.
		\end{itemize}
	\end{frame}
	
	% ---------------------------------------------------------------
	\begin{frame}{Question 92}
		Which technique explains an individual prediction by feature contribution?
		\begin{enumerate}[label=\Alph*.]
			\item Elbow plot
			\item SHAP values
			\item ROC curve
			\item Confusion matrix
		\end{enumerate}
	\end{frame}
	\begin{frame}{Answer 92}
		\textbf{Correct: B} --- SHAP attributes a prediction to features.\\
		\vspace{0.5em}
		\textbf{Wrong:}
		\begin{itemize}
			\item A: used for choosing k.
			\item C: summarizes threshold performance.
			\item D: summarizes classification outcomes.
		\end{itemize}
	\end{frame}
	
	% ---------------------------------------------------------------
	\begin{frame}{Question 93}
		A model must score 100 million Delta records overnight. Which approach is best?
		\begin{enumerate}[label=\Alph*.]
			\item \texttt{toPandas()} then loop on the driver
			\item Manual scoring in a spreadsheet
			\item Batch inference with Spark, for example \texttt{spark\_udf} or \texttt{fe.score\_batch()}
			\item A real-time REST endpoint per row
		\end{enumerate}
	\end{frame}
	\begin{frame}{Answer 93}
		\textbf{Correct: C} --- batch scoring is efficient for large offline jobs.\\
		\vspace{0.5em}
		\textbf{Wrong:}
		\begin{itemize}
			\item A: does not scale.
			\item B: not feasible.
			\item D: too slow and costly.
		\end{itemize}
	\end{frame}
	
	% ---------------------------------------------------------------
	\begin{frame}{Question 94}
		A team needs low-latency REST predictions. What should they create?
		\begin{enumerate}[label=\Alph*.]
			\item A Git folder
			\item A Model Serving endpoint
			\item A SQL dashboard
			\item A scheduled batch job
		\end{enumerate}
	\end{frame}
	\begin{frame}{Answer 94}
		\textbf{Correct: B} --- serving endpoints expose models as REST APIs.\\
		\vspace{0.5em}
		\textbf{Wrong:}
		\begin{itemize}
			\item A: not deployment.
			\item C: not a model API.
			\item D: not low latency.
		\end{itemize}
	\end{frame}
	
	% ---------------------------------------------------------------
	\begin{frame}{Question 95}
		A team needs continuous scoring as new data arrives. What is a suitable approach?
		\begin{enumerate}[label=\Alph*.]
			\item Spark Structured Streaming with a model UDF
			\item Exporting to CSV weekly
			\item A one-time \texttt{OPTIMIZE}
			\item Manually re-running a notebook hourly
		\end{enumerate}
	\end{frame}
	\begin{frame}{Answer 95}
		\textbf{Correct: A} --- streaming applies the model to incoming micro-batches.\\
		\vspace{0.5em}
		\textbf{Wrong:}
		\begin{itemize}
			\item B: not streaming.
			\item C: unrelated.
			\item D: not continuous.
		\end{itemize}
	\end{frame}
	
	% ---------------------------------------------------------------
	\begin{frame}{Question 96}
		Which request key is valid when sending data to a Databricks serving endpoint as JSON?
		\begin{enumerate}[label=\Alph*.]
			\item \texttt{model\_input\_df}
			\item \texttt{records\_pandas}
			\item \texttt{dataframe\_records}
			\item \texttt{spark\_rows}
		\end{enumerate}
	\end{frame}
	\begin{frame}{Answer 96}
		\textbf{Correct: C} --- other valid formats include \texttt{dataframe\_split} and \texttt{inputs}.\\
		\vspace{0.5em}
		\textbf{Wrong:}
		\begin{itemize}
			\item A: not valid.
			\item B: not valid.
			\item D: not valid.
		\end{itemize}
	\end{frame}
	
	% ---------------------------------------------------------------
	\begin{frame}{Question 97}
		A team wants to split traffic between two model versions on one endpoint. What should they configure?
		\begin{enumerate}[label=\Alph*.]
			\item Traffic percentages across served entities
			\item Two separate experiments
			\item Different Git branches
			\item Two feature tables
		\end{enumerate}
	\end{frame}
	\begin{frame}{Answer 97}
		\textbf{Correct: A} --- an endpoint can serve multiple versions with a traffic split.\\
		\vspace{0.5em}
		\textbf{Wrong:}
		\begin{itemize}
			\item B: unrelated.
			\item C: unrelated.
			\item D: unrelated.
		\end{itemize}
	\end{frame}
	
	% ---------------------------------------------------------------
	\begin{frame}{Question 98}
		What do inference tables capture for a serving endpoint?
		\begin{enumerate}[label=\Alph*.]
			\item Notebook revisions
			\item Request and response payloads logged to a Delta table
			\item Training data versions
			\item Cluster logs only
		\end{enumerate}
	\end{frame}
	\begin{frame}{Answer 98}
		\textbf{Correct: B} --- they support monitoring and debugging.\\
		\vspace{0.5em}
		\textbf{Wrong:}
		\begin{itemize}
			\item A: not their purpose.
			\item C: not their purpose.
			\item D: not their purpose.
		\end{itemize}
	\end{frame}
	
	% ---------------------------------------------------------------
	\begin{frame}{Question 99}
		What is the difference between data drift and concept drift?
		\begin{enumerate}[label=\Alph*.]
			\item Data drift is a change in labels only
			\item Concept drift only happens during training
			\item They are identical terms
			\item Data drift is a change in input distributions; concept drift is a change in the feature-target relationship
		\end{enumerate}
	\end{frame}
	\begin{frame}{Answer 99}
		\textbf{Correct: D} --- both can degrade performance in production.\\
		\vspace{0.5em}
		\textbf{Wrong:}
		\begin{itemize}
			\item A: it concerns inputs.
			\item B: it happens after deployment.
			\item C: they differ.
		\end{itemize}
	\end{frame}
	
	% ---------------------------------------------------------------
	\begin{frame}{Question 100}
		Which Databricks capability monitors tables and inference data for drift and quality?
		\begin{enumerate}[label=\Alph*.]
			\item Git folders
			\item \texttt{dbutils.fs}
			\item Cluster autoscaling
			\item Lakehouse Monitoring
		\end{enumerate}
	\end{frame}
	\begin{frame}{Answer 100}
		\textbf{Correct: D} --- it tracks statistics and drift over time.\\
		\vspace{0.5em}
		\textbf{Wrong:}
		\begin{itemize}
			\item A: version control.
			\item B: file utilities.
			\item C: compute scaling.
		\end{itemize}
	\end{frame}
	
\end{document}